\documentclass[12pt]{report}

% Packages nécessaires
\usepackage[utf8]{inputenc}
\usepackage[T1]{fontenc}
\usepackage[english]{babel}
\usepackage{titlesec}
\usepackage{tocloft}
\usepackage{lipsum} % Package de remplissage de texte (à retirer dans le rapport final)
\usepackage{graphicx}
\usepackage{float} % pour l'option de placement [H]
\usepackage[colorlinks=true, linkcolor=blue, citecolor=green]{hyperref}
\usepackage{minted} % pour la coloration syntaxique
\usepackage{wrapfig} % Pour l'enroulement du texte autour des images
\usepackage{array}
\usepackage{colortbl} %Pour les couleurs des colonnes et lignes des tableaux
% Définition de nouveaux types de colonnes centrées
\newcolumntype{C}[1]{>{\centering\arraybackslash}p{#1}}

% Mathématiques
\usepackage{amsmath, amssymb}
\usepackage{amsthm}
\usepackage{mdframed}  % Package pour les cadres

% Environnements pour les théorèmes, définitions et preuves
\newtheorem{theorem}{Théorème}
\newtheorem{definition}{Définition}
\newtheorem{example}{Exemple}
\newtheorem{property}{Propriété}

\newmdtheoremenv{boxedproperty}{Propriété}[section]  % Crée un environnement encadré pour les propriétés

% Configuration de la page
\usepackage[a4paper, left=2.5cm, right=2.5cm, top=2.5cm, bottom=2.5cm]{geometry}

% Pour avoir le petit dessin de l'INSA en bas à droite de chaque page
\usepackage{background}
\backgroundsetup{
scale=1,
angle=0,
opacity=1,
color=black,
contents={\begin{tikzpicture}[remember picture,overlay]
\node at ([xshift=-0.8in,yshift=0.8in] current page.south east) % Adjust the position of the logo.
{\includegraphics[scale=0.8]{images/pattern.png}}; % logo goes here
\end{tikzpicture}}
}

% Configuration des titres des sections et sous-sections
\titleformat{\chapter}[display]
  {\normalfont\huge\bfseries}{\chaptername\ \thechapter}{20pt}{\Huge}
\titleformat{\section}
  {\normalfont\Large\bfseries}{\thesection}{1em}{}
\titleformat{\subsection}
  {\normalfont\large\bfseries}{\thesubsection}{1em}{}

% Numérotation des chapitres en chiffres romains
\renewcommand{\thechapter}{\Roman{chapter}.}
% Numérotation des sections et sous-sections en chiffres arabes
\renewcommand{\thesection}{\arabic{section}.}
\renewcommand{\thesubsection}{\arabic{section}.\arabic{subsection}}

% Configuration de la table des matières
\renewcommand{\cftchapfont}{\bfseries}
\renewcommand{\cftsecfont}{\normalfont}
\renewcommand{\cftsubsecfont}{\normalfont}
\renewcommand{\cftchappagefont}{\bfseries}
\renewcommand{\cftsecpagefont}{\normalfont}
\renewcommand{\cftsubsecpagefont}{\normalfont}
\setlength{\cftbeforetoctitleskip}{0pt}
\setlength{\cftaftertoctitleskip}{10pt}
\renewcommand{\contentsname}{Table des matières}

\begin{document}
    % Page de titre
\begin{titlepage}
    \centering

    % Minipages pour les logos
    \noindent % Assure qu'il n'y a pas d'indentation au début de la ligne
    \begin{minipage}{0.5\textwidth}
        \includegraphics[width=0.5\linewidth]{images/logo_INSA.png} % Ajustez le chemin et la taille
    \end{minipage}%
    \hfill % Assure que les deux minipages seront poussés à l'extrême gauche et droite
    \begin{minipage}{0.5\textwidth}
        \flushright % Alignement à droite dans la minipage
        \includegraphics[width=0.5\linewidth]{images/logo_universite_toulouse.png} % Ajustez le chemin et la taille
    \end{minipage}


    \vspace*{2cm} % Espace vertical de 2 cm
    {\Huge\bfseries Innovative Smart System \par}
    \vspace{1cm}
    {\huge Portefolio\par}
    \vspace{2cm}
    {\Large \textbf{Clément Gauché} \par}
    \vspace{2cm}
    
    \begin{figure}[H]
    \centering
    \includegraphics[width=0.7\textwidth]{images/Logo_ISS_sans_texte.png}
    \end{figure}
    
    \vfill
    {\Large \textbf{Institut National des Sciences Appliquées de Toulouse} \par}
    \vspace{1cm}
    {\large \today \par}
\end{titlepage}

% Résumé concis, souvent utilisé pour donner un aperçu rapide du contenu d'une recherche ou d'un article académique.
\begin{abstract}
    \lipsum[1] % Lorem Ipsum
\end{abstract}

% Table des matières
\tableofcontents

% Introduction
\chapter{Introduction}
Ce portfolio est organisé de la façon suivante.\\
Chaque section correspond a une matière, et est décomposée en sous parties:
\begin{itemize}
    \item Context
    \item Technical summary
    \item Pratical Work
    \item Competences gained
    \item Analysis and remarks
    \item Reflections
\end{itemize}

% Début du contenu du rapport
\chapter{Courses description}
\section{Layout of the year's courses}
\begin{figure}[H]
    \centering
    \begin{tabular}{|C{2cm}|C{2cm}|C{4cm}|C{6cm}|}
        \hline
        \rowcolor{black!15!white}\textbf{Date} & \textbf{Duration} & \textbf{Context} & \textbf{Content} \\
        \hline
        from 15/06 to 22/09/2024 & 576 hours & Internship at IASI in Romania & Trainee in the network department, Implementation of monitoring functions in a Software Defined Network (SDN) \\
        \hline
        from 30/09/2024 to 29/01/2025 & 68.5 hours & Innovative Project & The guideline industrial project of our year \\
        \hline
        from 30/09 to 15/10/2024 & 21.5 hours & Cloud and Edge computing & Learn and practice about virtualization techniques such as VMs and containers \\
        \hline
        from 30/09 to 03/12/2024 & 7.5 hours & 5G Technologies & Learn about 5G and more generally topics related to cellular networks \\
        \hline
        from 04/10/2024 to 14/01/2025 & 33.5 hours & Service Architecture & Learn about legacy and modern architectures service oriented for software engineering \\
        \hline
        from 07/10 to 14/12/2024 & 18.25 hours & Wireless Sensor Networks & Learn about wireless sensor networks technologies \\
        \hline
        from 15/10 to 09/12/2024 & 14.75 hours & Middleware for IoT & Discover communication protocols for IoT \\
        \hline
        from 15/10 to 24/01/2025 & 34.5 hours & Security for connected objects & Discussion about the need and how to secure protocols for IoT \\
        \hline
        from 07/10 to 17/10/2024 & 10.5 hours & Energy for Connected Objects & Introduction to the different methods to power Energy for Connected Objects \\
        \hline
        
    \end{tabular}
\end{figure}


\begin{figure}[H]
    \centering
    \begin{tabular}{|C{2cm}|C{2cm}|C{4cm}|C{6cm}|}
        \hline
        \rowcolor{black!15!white}\textbf{Date} & \textbf{Duration} & \textbf{Context} & \textbf{Content} \\
        \hline
        from 04/11 to 08/11/2023 & 13.75 hours & Lab at AIME & Discover the creation of a sensor \\
        \hline
        from 12/11 to 18/12/2024 & 39 hours & Microcontrollers Open-Source Hardware and Sensor Introduction & Introduction to microcontrollers programming and implementation of our sensor in a complex circuit \\
        \hline
        from 04/12/2024 to 07/01/2025 & 7.5 hours & Communication protocols for LP-WPAN & Definition of a TCP/IP protocol stack for an LP-WPAN network \\
        \hline
        from 04/12/2024 to 22/01/2025 & 15.75 hours & Embedded IA for IoT & Introduction to the concepts of AI applied to an IoT context \\
        \hline
        from 08/01 to 24/01/2024 & 10.5 hours & Emerging network (SDN, NGN) & Discovery of emerging network paradigms \\
        \hline
        from 05/10 to 31/01/2024 & 6.75 hours & Portfolio & Writing of the document you are actually reading \\
        \hline
    \end{tabular}
    \caption{Table with information about all my projects}
\end{figure}


\newpage
% Cloud and Edge Computing
\section{Cloud and Edge Computing}
\subsection{Context}
%Le cours sur le Cloud and Edge Computing a été le premier cours que j'ai eu en 5ISS. Il introduit les concepts clés de l'informatique distribuée et omniprésente. Il met en lumière des architectures telles que le « Service Oriented Computing (SOC) », la virtualisation, ou encore le rôle des technologies cloud dans le développement de l'informatique omniprésente et utilitaire.
\indent \indent Among the first courses I came across during my final year at 5ISS was Cloud and Edge Computing. This course introduced me to core concepts of both distributed and pervasive computing together with the analysis of a basic architecture that included techniques of Service Oriented Computing (SOC), virtualization techniques, and cloud technologies as enablers of utility and ubiquitous computing. By nature, the course had both a theoretical and hands-on combination that gave an insight into how such paradigms shape the current IT systems and their applications.

\subsection{Technical Summary} % = synthèse du cours
%- Virtualization: Concepts, hypervisors (Type 1 and 2), paravirtualization, containerization, and VM templates. \\
%Comparison between "Service Orchestration" and "Service Choregraphy". -> Which is the best? - > It depends on the context: \\
%The orchestrator slows down the system but is easier to update.  \\
%- > Depends on the development model and on the economic model. \\
%- Cloud Computing: NIST and other definitions, core characteristics such as on-demand service, scalability, pay-per-use models, and service delivery (IaaS, PaaS, SaaS). \\
%- Edge Computing: Integration with IoT, the distribution of data control, benefits (e.g., latency reduction, data sovereignty), and challenges like connectivity and limited processing power. \\
%- Fog Computing: Bridging cloud services closer to data sources for improved real-time interaction and reduced data travel distance. "Cloudifying" the IOT world !!!

\indent \indent Virtualization has been amply covered in class. To be more specific, we discovered the different types of it: hypervisors, Type 1 and 2, paravirtualization, and containerization. We have also considered Virtual Machine Templates, which strongly ease virtual environments management and deployment. A specific highlight was the comparison between service orchestration and service choreography, where we analyzed the strengths and limitations of each. While orchestration does indeed ease updates by adopting a centralized approach, it tends to make the system slower. On the other hand, choreography is decentralized, enhancing autonomy but with a very painstaking design process. The choice between them is necessarily dependent on the specific contexts of development and economics in which they are applied.\\
\vspace{0.25cm}

Another important area discussed in this course was cloud computing: a discussion of definitions, among them the widely recognized framework from NIST, but also established its core characteristics. On-demand service provision, scalability, pay-per-use, and supply of services in models like IaaS, PaaS, and SaaS-all these form key elements that make cloud computing. We also looked at edge computing, which is integral to the IoT ecosystem, and allows for latency reduction and ensures data sovereignty. However, this also presents challenges of connectivity and the processing power of the edge devices. Finally, we discussed fog computing: an intermediary layer that connects cloud services with data sources, reducing the distance data needs to travel, and allowing real-time interaction. This was aptly referred to as the process of "cloudifying the IoT world."

\subsection{Practical work}
%La partie pratique s'est organisée en 2 TP, un sur la découverte des Cloud Hypervisors et l'autre sur l'orchestration de services dans un environnement hybride cloud/edge. Ces TP m'ont permis d'approfondir mes connaissances des concepts et des technologies utilisées dans les différentes techniques de virtualisation.
\indent \indent The practical work for this course was divided into two major sessions, each designed to reinforce the theoretical aspects covered in class. In the first session, we explored cloud hypervisors by engaging with tools such as VirtualBox and the OpenStack API. This practical session allowed me to experiment with various virtualization techniques and understand their implementation in controlled environments. The second session was on service orchestration in hybrid cloud and edge environments. In that exercise, I learned the deployment and management of a distributed system operating across these two architectures, thus providing insights into their complexities and potential.
\vspace{0.25cm}

These practical sessions were crucial for bridging the gaps in theory into applications. These have provided a chance not only to comprehend the conceptual aspects but also to apply the same in realistic scenarios. This solidifies my comprehension and confidence in handling Hybrid Architectures.

\subsection{Skills acquired}
%- Proficiency in cloud and edge computing architecture design. \\
%- Skills in using virtualization tools and managing virtual environments (VMs and containers). \\
%- Expertise in deploying services using container orchestration platforms like Kubernetes. \\
%- The ability to design resilient and low-latency systems tailored for IoT and distributed scenarios. \\
%- Insight into handling edge node variability, ensuring service continuity despite real-time disruptions.
\indent \indent The course really enhanced my technical capability and comprehension of Distributed Systems. Cloud and Edge computing architectures focusing on high performance and dependability. My experience with VirtualBox and OpenStack gave me hands-on experience in creating and managing virtual environments: virtual machines and containers. This also includes the development of expertise in deploying and orchestrating containerized services through the use of platforms like Kubernetes, which is critical in modern-day designing of distributed systems. The course enhanced my system design skills for low latency and high resilience, particularly in IoT and Distributed Environments. Apart from that, I learned real-time disruption management at the edge and continuity of service, despite variability in edge nodes.

\subsection{Analysis and remarks}
%Le format des TP est pour moi très intérressant car étant divisé en une partie théorique, permettant de fixer les concepts essentiels du cours en début de TP, accompagné d'une partie pratique, manipulation Virtual Box et l'API OpenStack, tout en permettant de mettre en application ces concepts de cours.
\indent \indent One of the strong points of this course is the structuring of the practical sessions, which did combine theory and practice quite nicely. Each session started with a theoretical introduction to reinforce the central concepts of the course; we then had to go through an implementation phase by putting these ideas into practice with tools such as VirtualBox and OpenStack APIs. In this way, both my understanding was solidly shaped and practical skills were gained; hence, learning has been both engaging and very effective. I found this format particularly enriching because it mirrors real-world workflows, where a deep understanding of concepts must complement practical expertise.

\subsection{Reflections}
\indent \indent This course gave me a better understanding of the interaction between cloud and edge computing systems and how they enable innovative, scalable, and efficient solutions. The smooth integration of theoretical and practical aspects has prepared me for the challenges of designing and managing distributed systems in professional life. In this learning process, I developed not only my technical skills but also the problem-solving approach needed to handle modern IT system complexities.

\newpage
% 5G Technologies
\section{5G technologies}
\subsection{Context}

\indent \indent The 5G technologies module presented by Professor Etienne Sicard, focused on the rapid evolution of cellular networks, emphasizing the transformative nature of 5G. A unique aspect of the learning approach was the use of reverse pedagogy, where students presented the majority of the content. The central goal was to explore cutting-edge topics in mobile communications, ranging from modulation techniques to the societal impacts of 5G and beyond.

5G represents a paradigm shift in mobile networks, introducing Software Defined Radio (SDR) and microservices, which revolutionized network architecture by increasing flexibility and efficiency. Through this course, we students, were encouraged to examine the current state of cellular technology while considering its implications for the future.
\subsection{Technical and practical work}

\indent \indent During this course we could discuss about many varied subject. I collaborated with Noel Jumin on a presentation titled "Samsung's Vision for 6G," analyzing how Samsung's ambitions for 6G diverged from its competitors. We highlighted Samsung's focus on sub-THz frequencies, AI integration, and next-generation services such as holographic communications. This required an in-depth understanding of how 5G technologies paved the way for these future advancements, including the challenges posed by scaling up.

Other students' presentations covered topics such as 5G modulation techniques, vehicular networks, and the environmental and societal implications of 5G and 6G. These presentations provided a comprehensive understanding of both the technical and non-technical aspects of mobile communications.

\subsection{Skills acquired}
\subsection{Analysis and remarks}
This course was highly engaging as it dealt with contemporary and practical topics in cellular network technology. The discussion of 5G as a foundational step towards 6G was particularly fascinating, showing how each generation builds upon the previous one. The reverse pedagogy approach fostered a dynamic and interactive learning environment.
\\
However, I found reverse pedagogy less effective overall. While preparing and delivering our own presentations allowed us to master our chosen topic, it made it harder to fully comprehend the topics presented by others. During the presentations, my focus was often divided between understanding the other groups' content and refining my own work.
\\
I believe that a better balance between lectures by the teacher and student presentations would enhance the learning experience. Teacher-led lectures could provide a deeper dive into each topic, ensuring a more comprehensive understanding, while our presentations could supplement this with detailed case studies or specific insights. As it stands, I feel I gained high-level knowledge of the subjects presented by others, but lacked the depth that direct teaching could have provided.

\newpage
% Service orienté architecture
\section{Service Architecture - Software Engineering}
\subsection{Context}
3 courses : \\
The SOAP standard: In this course, you have studied the SOA architecture, the web service technology and its standards. \\
RESTFull services: \\
Microservices: How can monolithic applications can be limiting in developping flexible and scalable applications. I have studied SOA and REST architecture and what they bring to overcome the monolithic issues. Finally, I have discovered the Microservices architectures that extends the SOA architecture.
\subsection{Technical = synthèse du cours}
\subsection{Practical work = Projets de TP}
\subsection{Skills acquired}
\subsection{Analysis and remarks}

\newpage
% Wireless Sensor Network
\section{Wireless Sensor Networks}
\subsection{Context}
\subsection{Technical}
\subsection{Practical work}
\subsection{Skills acquired}
\subsection{Analysis and remarks}

\newpage
% Middleware for IOT
% \section{Middleware for IoT}
% \subsection{Context}
% Le cours de Middleware for IoT a été donné par le professeur Thierry Monteil. Les objectifs principaux de ce cours sont de découvrir différentes approches et protocoles de communication, et de comprendre les enjeux de la communication entre objets connectés afin d'être capable de mettre en place des solutions adaptées au besoin.\\
% Ce cours, comme pour le cours de \hyperref[sec:service_architecture]{Service Architecture} donné par Ms. Nawal Guermouche, était sous la forme d'un MOOC (Massive Open Online Course), à la suite duquel nous avons eu une serie de travaux pratiques. Ces TP qui visent a illustrer les concepts abordés dans le MOOC, notamment l'utilisation d'MQTT, oneM2M et de Node-RED.\\
% L'objectif des TP était de comprendre comment faire intéragir plusieurs objets connectés entre eux, en choisisant les protocoles adaptés, tout en prenant en compte la diversité des infrastructures matérielles et logicielles exixtantes.

% \subsection{Technical Summary}
% Comme dit précédemment, durant ce cours nous avons abordés plusieurs protocoles de communication, et notamment Message Queuing Telemetry Transport (MQTT). Il s'agit d'un protocole de messagerie léger, conçu pour les appareils ayant des ressources limitées. Son fonctionnement est basé sur un broker, un serveur, qui gère des sujets de publication (des topics). Les clients "publishers" publient des messages sur des topics, et les clients "subscribers" recoivent des messages en s'abonnant à des topics.\\
% Grace à son architecture simple, le protocole MQTT est très utilisé dans l'IoT, où la consommation de bande passante doit être minimisée, et la fiabilité y être essentielle.\\

% Le second protocole que nous avons étudié est oneM2M. Il s'agit plus précisemment d'une norme soutenue par plusieurs organismes de standardisation, qui vise à fournir et assurer l'intéropérabilité des architectures IoT. Basée sur REST, oneM2M propose une structure hiérarchique permettant de modéliser des ressources pour une application IoT. Les ressources sont organisées en trois niveaux : les Application Entities (AE), les Containers (CNT) et les Content Instances (CI). Ce sont ces niveaux définissent l'espace de stockage hiérarchique des Servers. Les Application Entities (AE) représentent les fonctionnalités disponibles dans lesquelles il y a des containers (CNT) qui sont des catégories de données et des Content Instance (CI) qui sont des instances de données publiée.\\
% Ainsi, c'est au travers d'un container dans une AE, qu'un appareil IoT peut publier ou lire des informations.\\

% \subsection{Practical Work}
% Le premier projet des séances de TP, a été le déploiemment d'un broker MQTT, avec le logiciel Mosquitto, et la mise en place d'échanges de messages simple entre plusieurs ESP8266. Un ESP8266 jouant le role du bouton, le publisher, et un autre controllant une LED, le subscriber. Chaque changement de l'état du bouton était reflété sur l'extinction ou l'allumage de la LED.\\

% Le deuxième projet, qui était sensiblement le même que le premier mais avec oneM2M, a commencé par le suivit d'un tuto sous un Jupyter Notebook, pour en comprendre les concepts. Les premiers exercices de ce notebook nous ont permis de comprendre comment créer des AE, des containers, et des content instances. Puis, nous sommes allé plus loin en découvrant les créations de groupes, les notifications et les règles qui permettenht de mettre a jour les ressources de manière automatique.\\
% Ce n'est qu'une fois ces concepts  maîtrisés que nous avons utilisé un serveur oneM2M et un script Python pour simuler la même application qu'avec MQTT. Afin d'avoir quelque chose d'équivalent au premier projet, mon binome et moi avons simulé un “bouton” et une “LED” virtuelle. Un script Python nous permettait de créer des AE, des containers, et des content instances pour le bouton et la LED, et de gérer les notifications. 

% Si nous avions eu plus de temps sur ce projet, j'aurais aimé pouvoir intégrer Node-RED pour pouvoir manager plus facilement les flux de données entre les différents appareils IoT. Node-RED nous aurait apporté une interface visuelle, connectant les dispositifs matériels, les API et des services en ligne de manière intuitive.

% \subsection{Skills Acquired}
% \begin{table}[H]
%     \centering
%     \begin{tabular}{|p{3.5cm}|p{8cm}|p{3.5cm}|}
%     \hline
%     \textbf{Skill} & \textbf{Description} & \textbf{Level of Mastery} \\
%     \hline
%     Situate IoT Standards & Know how to situate the main standards for the Internet of Things & Advanced \\
%     \hline
%     Deploy IoT Architecture & Deploy an architecture compliant to an IoT standard and implement a sensor network & Intermediate \\
%     \hline
%     Configure OM2M & Deploy and configure an IoT architecture using OM2M & Intermediate \\
%     \hline
%     Interact with REST & Interact with the different resources of the architecture using REST services & Advanced \\
%     \hline
%     Integrate New Technology & Integrate a new technology into the deployed architecture & Intermediate \\
%     \hline
%     \end{tabular}
%     \caption{Skills Acquired}
%     \label{tab:skills_acquired}
% \end{table}


% \subsection{Analysis and remarks}

% La découverte et l'expérimentation de deux protocoles IoT complémentaires, MQTT (léger, adapté à des cas d'usage rapide et peu gourmands en ressources) et oneM2M (plus structuré et standardisé, offrant un cadre complet pour des applications IoT à grande échelle), m'ont permis de mieux saisir l'importance de l'interopérabilité dans l'Internet des Objets. La mise en pratique a été particulièrement instructive, de la configuration d'un broker Mosquitto, à la conception d'un mini-projet sur des ESP8266 et l'exploration de oneM2M au travers de scripts Python, en passant par Node-RED pour la mise en flux des données.
% Toutefois, j'aurais aimé approfondir davantage l'intégration directe de oneM2M dans Node-RED pour mieux en exploiter tout le potentiel, notamment en termes de notifications et de règles de gestion.

% Par ailleurs, l'installation des outils sur les machines personnelles nous a freinés, mon binome et moi, notamment sur la partie tuto oneM2M sur le notebook Jupyter. J'aurai aimé qu'une solution fonctionnelle, telle qu'une machine virtuelle à déployer, nous soit fourni en début de TP comme cela est fait notamment dans les TP du \hyperref[sec:service_architecture]{Service Architecture}.

% Sur le plan personnel, cette matière vient enrichir mon parcours en me donnant une vue plus globale des différentes approches de communication dans l'IoT. Les compétences acquises me seront utiles pour de futurs projets où il faudra choisir la solution la plus adaptée à un contexte précis.

% Souhaitant m'orienter vers le Master Spécialisé en TLS-Sec à l'issue de mon diplôme INSA, je regrette que la dimension sécurité de ces protocoles n'aient pas été abordé durant les travaux pratiques, car la sécurisation des échanges (chiffrement, authentification, protection contre les intrusions) constitue un enjeu majeur en IoT.

\section{Middleware for IoT}
\subsection{Context}
The Middleware for IoT course was taught by Professor Thierry Monteil. The main objectives of this course were to explore different approaches and communication protocols, and to understand the challenges of communication between connected devices in order to implement solutions tailored to specific needs.

Like the course \hyperref[sec:service_architecture]{Service Architecture} taught by Ms.\ Nawal Guermouche, this course was delivered in the form of a MOOC (Massive Open Online Course). Afterwards, we carried out a series of practical exercises designed to illustrate the concepts discussed in the MOOC, notably the use of MQTT, oneM2M, and Node-RED.

The goal of these practical sessions was to understand how multiple connected devices can interact with each other by choosing appropriate protocols, while taking into account the variety of existing hardware and software infrastructures.

\subsection{Technical Summary}
As mentioned earlier, we studied several communication protocols during this course, particularly Message Queuing Telemetry Transport (MQTT). This is a lightweight messaging protocol designed for devices with limited resources. Its functioning relies on a broker (a server) that manages publication topics. Client publishers publish messages on specific topics, and client subscribers receive messages by subscribing to these topics.

Thanks to its simple architecture, the MQTT protocol is widely used in IoT, where bandwidth consumption must be minimized and reliability is essential.

The second protocol we studied is oneM2M. More precisely, it is a standard endorsed by multiple standardization bodies, aiming to provide and ensure interoperability among IoT architectures. Based on REST, oneM2M proposes a hierarchical structure to model resources for an IoT application. The resources are organized on three levels: Application Entities (AE), Containers (CNT), and Content Instances (CI). These levels define the hierarchical storage space on servers. Application Entities (AE) represent available functionalities, within which there are Containers (CNT) that serve as data categories and Content Instances (CI) that hold published data instances.

Thus, it is through a container in an AE that an IoT device can publish or read information.

\subsection{Practical Work}
The first project in the practical sessions involved deploying an MQTT broker using Mosquitto and setting up simple message exchanges among several ESP8266 devices. One ESP8266 acted as a button (the publisher), while another controlled an LED (the subscriber). Any change in the button’s state was reflected by turning the LED on or off.\\

The second project, which was essentially the same as the first but with oneM2M, began with following a tutorial in a Jupyter Notebook to understand the core concepts. The first exercises in the notebook enabled us to learn how to create AEs, containers, and content instances. We then went a step further by exploring group creation, notifications, and rules that can automatically update resources.\\

Only after mastering these concepts did we use a oneM2M server and a Python script to simulate the same application we had built with MQTT. To achieve an equivalent setup to the first project, my partner and I simulated a virtual “button” and a virtual “LED.” A Python script allowed us to create AEs, containers, and content instances for both the button and the LED, as well as manage notifications. 

Had we had more time for this project, I would have liked to integrate Node-RED to more easily manage data flows among the various IoT devices. Node-RED would have provided a visual interface to intuitively connect hardware devices, APIs, and online services.

\subsection{Skills Acquired}

\begin{table}[h!]
    \centering
    \arrayrulecolor{black} % Defines the color of the table borders
    \renewcommand{\arraystretch}{1.5} % Adjusts the vertical spacing between rows
    \begin{tabular}{|p{11cm}|c|c|}
    \hline
    \rowcolor[gray]{0.8}
    \textbf{Skills} & \textbf{Required} & \textbf{Achieved} \\ \hline
    Know how to situate the main standards for the Internet of Things & 4 & 4 \\ \hline
    Deploy an architecture compliant to an IoT standard and implement a sensor network & 4 & 4 \\ \hline
    Deploy and configure an IoT architecture using OM2M & 4 & 4 \\ \hline
    Interact with the different resources of the architecture using REST services & 4 & 4 \\ \hline
    Integrate a new technology into the deployed architecture & 4 & 4 \\ \hline
    \end{tabular}
    \caption{Skill matrix for Middleware for the Internet of Things}
    \label{table:skills-middleware-iot}
\end{table}

\begin{table}[H]
    \centering
    \arrayrulecolor{black} % Defines the color of the table borders
    \renewcommand{\arraystretch}{1.5} % Adjusts the vertical spacing between rows
    \begin{tabular}{|>{\raggedright}p{3.5cm}|p{8cm}|c|}
    \hline
    \rowcolor[gray]{0.8}
    \textbf{Skill} & \textbf{Description} & \textbf{Level of Mastery} \\
    \hline
    Situate IoT Standards & Know how to position the main standards for the Internet of Things & Advanced \\
    \hline
    Deploy IoT Architecture & Deploy an architecture compliant with an IoT standard and implement a sensor network & Intermediate \\
    \hline
    Configure OM2M & Deploy and configure an IoT architecture using OM2M & Intermediate \\
    \hline
    Interact with REST & Interact with different resources of the architecture using REST services & Advanced \\
    \hline
    Integrate New Technology & Integrate a new technology into the deployed architecture & Intermediate \\
    \hline
    \end{tabular}
    \caption{Skills Acquired}
    \label{tab:skills_acquired}
\end{table}

\subsection{Analysis and Remarks}
Discovering and experimenting with two complementary IoT protocols—MQTT (lightweight, suitable for quick use cases with minimal resource consumption) and oneM2M (more structured and standardized, offering a comprehensive framework for large-scale IoT applications)—has helped me understand the importance of interoperability in the Internet of Things. The practical work was especially instructive, from configuring a Mosquitto broker and designing a mini-project on ESP8266 devices, to exploring oneM2M via Python scripts and using Node-RED for data flow management.

However, I would have liked to delve deeper into the direct integration of oneM2M into Node-RED to fully leverage its potential, particularly in terms of notifications and management rules.

In addition, installing the necessary tools on our personal machines slowed down my partner and me, especially for the oneM2M tutorial in the Jupyter Notebook. I would have appreciated having a ready-to-use solution—such as a virtual machine—to deploy at the start of the practical sessions, similar to what was provided for the \hyperref[sec:service_architecture]{Service Architecture} TP sessions.

From a personal standpoint, this course enhanced my academic path by giving me a more comprehensive view of different communication approaches in IoT. The skills I acquired will be valuable for future projects that require choosing the most appropriate solution for a specific context.

Intending to pursue the Master’s Program in TLS-Sec after completing my INSA degree, I regret that the security aspects of these protocols were not addressed during the practical work. Ensuring secure exchanges (encryption, authentication, intrusion protection) is a major challenge in IoT.


\newpage
% Security for Connected Objects
\section{Security for connected objects}
\subsection{Context}
\subsection{Technical}
\subsection{Practical work}
\subsection{Skills acquired}
\begin{table}[h!]
    \centering
    \arrayrulecolor{black} % Defines the color of the table borders
    \renewcommand{\arraystretch}{1.5} % Adjusts the vertical spacing between rows
    \begin{tabular}{|p{11cm}|c|c|}
    \hline
    \rowcolor[gray]{0.8}
    \textbf{Skills} & \textbf{Required} & \textbf{Achieved} \\ \hline
    \rowcolor[gray]{0.9} \textbf{Security for connected objects} &  &  \\ \hline
    Knowing the main issues in security for IoT & 3 & 3 \\ \hline
    Understand the terminology of security & 2 & 2 \\ \hline
    Being able to have a critical look at the design of a system from a security point of view & 2 & 2 \\ \hline
    Being able to understand a scientific article that explains a weakness or a security solution and to explain it & 2 & 2 \\ \hline
    \end{tabular}
    \caption{Skill matrix for Security for Connected Objects}
    \label{table:skills-security-iot}
\end{table}

\subsection{Analysis and remarks}

\newpage
% Energy for IOT (or Connected Objects)
\section{Energy for connected objects}
\subsection{Context}
% Le cours « Énergies pour Objets Connectés » a été donné par Gaël Loubet. Par le biais de plusieurs séances théoriques et de deux travaux pratiques, nous avons étudié les différentes solutions existantes pour alimenter des dispositifs IoT sans batterie ni câblage direct. On s'est notamment concentré sur la récupération d’énergie ambiante, au transfert d’énergie sans fil et aux moyens de rendre les systèmes véritablement autonomes en énergie.
% En parallèle, la question de la consommation énergétique a aussi été abordée dans le cours sur les réseaux de capteurs sans fil, qui insistait sur le rôle central de la gestion de l’énergie dans le choix et l’optimisation des protocoles de communication.

\indent \indent The “Energies for Connected Objects” course was given by Gaël Loubet. In this course, we explored possible existing solutions for powering IoT devices without batteries or wired systems. Once the basics had been laid down, the course focused on ambient energy harvesting, wireless energy transfer and various solutions for making systems energy autonomous.

This course was given in parallel with the course on wireless sensor networks, in which the issue of energy consumption was also highlighted as a decisive factor in the choice of protocol definition.

\subsection{Technical Summary}
% Au début, nous avons passé en revue les bases de l’électronique basse consommation, comme l’utilisation de condensateurs ou de supercondensateurs pour stocker l’énergie récupérée, ainsi que les différents mécanismes de transfert d’énergie sans fil, allant du transfert en champ proche (couplage capacitif et inductif) à un mode radiatif lointain via une rectenna. Nous nous sommes aussi intéressés aux techniques de récupération d’énergie ambiante (lumineuse, mécanique, thermique, ondes électromagnétiques) et aux stratégies d’optimisation. Nous avons notamment abordé la conception d’antennes adaptées à une fréquence précise ou d’optimisations logicielles pour réduire la consommation dans les objets connectés.
% Par ailleurs, nous avons vu que dans un contexte d’IoT, il est crucial de dimensionner et d’optimiser la gestion de l’énergie, car les sources ambiantes sont souvent fluctuantes et de faible densité énergétique.
%La récupération d’énergie étant intermittente, nous avons vu qu'un système tampon (généralement un condensateur) est nécessaire pour lisser l’alimentation et éviter les chutes de tension subites.

\paragraph{}It contained discussions on capacitors and supercapacators to buffer energy, wireless power transfer that has been highly evolved with a near-field mechanism by means of capacitive and inductive coupling to a far-field radiative by the use of a rectenna, and energy harvesting techniques using ambient energy through conversion of energy due to light, mechanical motion, thermal gradient, and electromagnetic waves. The course also integrated optimization strategies, such as antenna design to match the frequency and software optimizations that would reduce energy consumption in connected devices.

\subsection{Practical work}
% Les séances de TP ont consisté à concevoir et tester des systèmes énergétiques en laboratoire. Un exercice en particulier portait sur l’alimentation d’une LED via l’énergie RF ambiante, sans batterie et uniquement grâce à une antenne.
% Nous avons notamment fait des calculs de la consommation requise par la LED pour fonctionner à 25, 50, et 100% de sa luminosité max. Nous avons notamment découvert qu'elle consommait au maximum 40 mW, ce qui est comparable à la consommation d’un microcontrôleur tel que le ... ou le ... A partir de cette puissance, nous avons pu déterminer la capacité de stockage nécessaire pour alimenter la LED en continu.
% Nous avons ensuite ensuite analysé un circuit redresseur, choisi des antennes adaptées, et optimisé le transfert d’énergie pour un fonctionnement efficace. Nous avons utilisé GNURadio et des cartes de test pour mesurer la puissance RF réelle émise. La stratégie « stocker puis utiliser », réalisée avec un gestionnaire de puissance bq25504 et un convertisseur DC-DC TPS63031, a permis un tamponnage et une utilisation efficace de l’énergie dans des scénarios de faible puissance.

\paragraph{}The practical work for this course involved the design and testing of energy systems in the laboratory. One of the projects was to power an LED from ambient RF energy. This involved analyzing a rectifier circuit, choosing appropriate antennas, and optimizing the energy transfer for efficient operation. I used GNURadio and spectrum analyzers to measure RF power output, find the optimal operating frequencies, and test various energy storage configurations. The "store then use" strategy, as realized with a bq25504 power management unit and a TPS63031 DC-DC converter, allowed for efficient energy buffering and utilization in low-power scenarios.

\subsection{Skills acquired}
\paragraph{}In this course, I developed high-level technical expertise in low-power circuit design, energy harvesting systems, and wireless power transfer technologies. I acquired hands-on experience with laboratory tools and methodologies such as frequency sweeps, impedance matching, and antenna characterization. Beyond the technical skills, I learned to interpret complex system data and optimize designs for practical applications. The course also fostered innovative problem-solving skills, particularly in balancing theoretical frameworks with real-world constraints.

\begin{table}[H]
    \centering
    \arrayrulecolor{black} % Defines the color of the table borders
    \renewcommand{\arraystretch}{1.5} % Adjusts the vertical spacing between rows
    \begin{tabular}{|p{11cm}|c|c|}
    \hline
    \rowcolor[gray]{0.8}
    \textbf{Skills} & \textbf{Required} & \textbf{Achieved} \\ \hline
    \rowcolor[gray]{0.9} \textbf{Energy for Connected Objects} &  &  \\ \hline
    Understand and master optimisation of communication protocols for IoT with respect to energy limitations & 4 & 4 \\ \hline
    Mastering the architecture of an energy management system, simple storage, energy recovery, know how to size the storage element according to the specifications & 4 & 4 \\ \hline
    \end{tabular}
    \caption{Skill matrix for Protocols and Communication}
    \label{table:skills-iot-energy-management}
\end{table}

\subsection{Analysis and remarks}
\paragraph{}This class demonstrated the potentials and pitfalls of designing energy-autonomous IoT devices. Even though I implemented various systems, such as the energy harvesting LED with ease, environmental inputs had high variability that posed the biggest challenges. For instance, supplying power continuously in constantly fluctuating electromagnetic environments forced me to devise creative solutions but also highlighted the necessity of using hybrid energy systems. Furthermore, a central design challenge was determining the balance between optimizing the capture of energy and reliability of the system.

\subsection{Reflections}
\paragraph{}"Energy for Connected Objects" provided a great platform for innovative energy solution development within the IoT domain. It identified the key challenges arising from future technologies in relevance to sustainability and efficiency; therefore, this was focused on the latest energy system approaches and applications. In the future, my skills and knowledge will guide my work in energy-autonomous devices, mainly in scalable wireless sensor networks, but also including advanced material development for energy conversion and storage. These experiences crystallized my aspiration to further the frontier of IoT energy systems toward sustainable technological development.

\newpage
% LAB at AIME
\section{LAB at AIME}

\newpage
% Microcontrollers Open-Source Hardware and Sensor Introduction
\section{Microcontrollers Open-Source Hardware and Sensor Introduction}

\newpage
% Communication protocols for LP-WPAN
\section{Communication protocols for LP-WPAN}

\newpage
% Embedded IA for IOT
% \section{Embedded IA for IOT}
% \subsection{Context} % Ajouter que j'avais déjà suivi un cours sur LinkedIn Learning de la prof de Lyon.
% \indent \indent Le module Embedded IA for IOT, enseigné par le professeur Phillipe Leleux, est la suite d'une première introduction au machine learning que nous avons eu en 4ème année.
% Le cours s'est concentré cette année sur l'intégration de l'intelligence artificielle (IA) dans les systèmes embarqués.
% Le cours visait à fournir aux étudiants une compréhension globale des défis et des opportunités présentés par la combinaison de l'IA et de l'Internet des objets (IoT).
% L'objectif principal était d'explorer comment l'IA peut être utilisée pour améliorer les capacités des systèmes embarqués, leur permettant d'effectuer des tâches complexes de manière autonome.

% \subsection{Technical Summary} %= Résumé du cours
% \indent \indent Dans ce cours, nous avons abordés des concepts théoriques et pratiques essentiels à l'intégration de l'IA dans des dispositifs contraints.
% Nous avons notamment abordé le prétraitement des données, l'étude des caractéristiques. Ce qui inclu les techniques de réduction de la dimensionnalité, telles que l'analyse en composantes principales (PCA) ou la sélection de caractéristiques.
% \\
% Nous avons aussi abordé les méthodes d'analyse et de classification des données, avec notamment les algorithmes Dynamic Time Warping (DTW) et les réseaux de neurones récurrents (RNN).
% \\
% Nous avons abordé l'apprentissage profond, en particulier les réseaux de neurones convolutifs (CNN), qui sont adaptés aux environnements ayant des ressources limitées.
% \\
% Nous avons discuté sur les paradigmes d'apprentissage distribué, permettant l'apprentissage de modèles d'IA sur des appareils tout en préservant la confidentialité des données.
% Il a mis en évidence les compromis entre l'IA basée sur le nuage, l'IA en périphérie et l'IA sur appareil en termes de bande passante, de latence et de fiabilité.

% \subsection{Practical work} %= Travaux pratiques : à reformuler
% \indent \indent Nous avons eu 2 sessions de travaux pratiques pour ce cours. Le premier TP avait pour sujet la reconaissance et la classification d'activités humaines.
% Les données étaient colléctées par 3 capteurs, qui sont des accéléromètres, positionnés au poignet, au à la poitrine et à la cheville d'un individu.
% Grace à ces données, nous devions classer des activités : s'asseoir, se tenir debout, marcher, faire du vélo, de la marche nordique, passer l'aspirateur et repasser.
% Durant cette session nous avons pu mettre en application le Dynamic Time Warping (DTW), utilisé pour comparer et aligner les données de séries temporelle.
% Puis nous avons classé les activités a l'aide des K plus proches voisins (KNN).
% Nous avons aussi appliqué de la reduction de dimensionnalité avec l'ACP. Nous avons enfin mis en oeuvre un reseau neuronal (perceptron multicouche) pour la classification des activités.
% Les outils que nous avons utilisés dans ce TP sont Python, scikit-learn et numpy.

% Le sujet du deuxième TP était de la detection de chute à partir d'images, en se concentrant sur les techniques d'apprentissage profond.
% L'ensemble des données était un ensemble d'images capturants différentes situations de chute.
% Nous avons construit et entrainé un réseau neuronal convolutif (CNN) pour classer les images en 2 catégories : chute et non chute.
% Nous avons pu appliquer des techniques de transfert d'apprentissage et d'optimisation du modèle comme la quantification et l'élaguage, afin d'améliorer l'éfficacité de notre modèle (objectif de contraintes de puissance : des petits microcontrolleurs = systèmes embarqués).
% Les outils que nous avons découverts dans ce TP sont TensorFlow, Keras et TensorFlow Lite.

% \subsection{Skills acquired} %= Compétences acquises

% \begin{table}[h!]
%     \centering
%     \arrayrulecolor{black} % Defines the color of the table borders
%     \renewcommand{\arraystretch}{1.5} % Adjusts the vertical spacing between rows
%     \begin{tabular}{|p{11cm}|c|c|}
%     \hline
%     \rowcolor[gray]{0.8}
%     \textbf{Skills} & \textbf{Required} & \textbf{Achieved} \\ \hline
%     \rowcolor[gray]{0.9} \textbf{Embedded IA for IoT} &  &  \\ \hline
%     Understand the characteristics of supervised and unsupervised learning problems & 4 & 4 \\ \hline
%     Understand the main basic methods and algorithms to deal with these problems & 4 & 4 \\ \hline
%     Understand the specificities of AI at the edge & 4 & 4 \\ \hline
%     Understand the main optimization methods enabling the embedding of AI algorithms & 4 & 4 \\ \hline
%     Be able to use these methods through Python libraries to solve practical problems with IoT data & 4 & 4 \\ \hline
%     \end{tabular}
%     \caption{Skills matrix for various software and data-related training units}
%     \label{table:skills-embedded-ia}
% \end{table}

% \begin{table}[h!]
%     \centering
%     \arrayrulecolor{black} % Defines the color of the table borders
%     \renewcommand{\arraystretch}{1.5} % Adjusts the vertical spacing between rows
%     \begin{tabular}{|p{3.5cm}|p{8cm}|c|}
%     \hline
%         \rowcolor[gray]{0.8}
%         \textbf{Competence} & \textbf{Description} & \textbf{Level of Mastery} \\ \hline
%         Data Preprocessing & Cleaning, scaling, and transforming IoT data for machine learning tasks. & Advanced \\ \hline
%         Time-Series Analysis & Classification and alignment techniques for temporal data from IoT sensors. & Intermediate \\ \hline
%         Deep Learning for IoT & Designing and optimizing CNNs for resource-constrained environments. & Advanced \\ \hline
%         Edge AI and Federated Learning & Implementing distributed learning paradigms while preserving privacy and improving performance. & Intermediate \\ \hline
%         TinyML and Model Optimization & Applying pruning, quantization, and transfer learning to create lightweight AI models. & Advanced \\ \hline
%         Python Libraries & Proficiency with TensorFlow, TensorFlow Lite, scikit-learn, and Keras for embedded AI development. & Advanced \\ \hline
%     \end{tabular}
%     \caption{Skills Acquired During the Embedded AI for IoT Module}
% \end{table}

% \subsection{Analysis and remarks} %= Analyse et remarques
% Ce cours m'a beaucoup plu car il m'a permis de mieux comprendre certains des concept que j'avais découvert au cours d'une formation LinkedIn Learning réalisé par la professeur Céline Robardet de l'INSA Lyon.
% Les travaux pratiques, en particulier le projet de detection de chute, 
% Il était gratifiant de voir comment des techniques telles que l'optimisation CNN et l'apprentissage fédéré pouvaient être appliquées pour résoudre des problèmes IoT du monde réel.

\section{Embedded AI for IoT}
\subsection{Context}
\indent \indent The module "Embedded AI for IoT," taught by Professor Philippe Leleux, builds upon an initial introduction to machine learning that we took during the fourth year. 
This year, the course focused on integrating artificial intelligence (AI) into embedded systems. 
The aim was to provide students with a comprehensive understanding of the challenges and opportunities presented by the combination of AI and the Internet of Things (IoT).
The primary objective was to explore how AI can enhance the capabilities of embedded systems, enabling them to perform complex tasks autonomously.

\subsection{Technical Summary}
\indent \indent In this course, we covered essential theoretical and practical concepts related to integrating AI into constrained devices. 
We studied data preprocessing and feature engineering, which included dimensionality reduction techniques such as Principal Component Analysis (PCA) and feature selection. \\
We also explored methods for data analysis and classification, focusing on algorithms like Dynamic Time Warping (DTW) and Recurrent Neural Networks (RNNs). \\
The course delved into deep learning, particularly Convolutional Neural Networks (CNNs), which are well-suited for resource-limited environments. \\
We discussed distributed learning paradigms, enabling AI models to be trained on devices while preserving data privacy. 
The trade-offs between cloud-based AI, edge AI, and on-device AI in terms of bandwidth, latency, and reliability were also highlighted.

\subsection{Practical Work}
\indent \indent We had two practical sessions for this course. The first session focused on recognizing and classifying human activities. 
The data was collected from three sensors (accelerometers) placed on the wrist, chest, and ankle of an individual. 
Using this data, we classified activities such as sitting, standing, walking, cycling, Nordic walking, vacuuming, and ironing.
During this session, we implemented Dynamic Time Warping (DTW) to compare and align time-series data. 
We then classified the activities using k-Nearest Neighbors (k-NN) and applied dimensionality reduction with PCA. 
Finally, we implemented a neural network (multi-layer perceptron) for activity classification. 
The tools used in this session included Python, scikit-learn, and numpy.

The second practical session focused on fall detection using image data, emphasizing deep learning techniques. 
The dataset consisted of images capturing various fall scenarios. 
We built and trained a Convolutional Neural Network (CNN) to classify images into two categories: "fall" and "no fall." 
We applied transfer learning and model optimization techniques, such as quantization and pruning, to improve the model's efficiency (a key consideration for constrained embedded systems such as microcontrollers). 
The tools explored in this session included TensorFlow, Keras, and TensorFlow Lite.

\subsection{Skills Acquired}
\begin{table}[H]
    \centering
    \arrayrulecolor{black}
    \renewcommand{\arraystretch}{1.5}
    \begin{tabular}{|p{11cm}|c|c|}
    \hline
    \rowcolor[gray]{0.8}
    \textbf{Skills} & \textbf{Required} & \textbf{Achieved} \\ \hline
    \rowcolor[gray]{0.9} \textbf{Embedded AI for IoT} &  &  \\ \hline
    Understand the characteristics of supervised and unsupervised learning problems & 4 & 4 \\ \hline
    Understand the main basic methods and algorithms to deal with these problems & 4 & 4 \\ \hline
    Understand the specificities of AI at the edge & 4 & 4 \\ \hline
    Understand the main optimization methods enabling the embedding of AI algorithms & 4 & 4 \\ \hline
    Be able to use these methods through Python libraries to solve practical problems with IoT data & 4 & 4 \\ \hline
    \end{tabular}
    \caption{Skills matrix for various software and data-related training units}
    \label{table:skills-embedded-ia}
\end{table}

\begin{table}[H]
    \centering
    \arrayrulecolor{black}
    \renewcommand{\arraystretch}{1.5}
    \begin{tabular}{|p{3.5cm}|p{8cm}|c|}
    \hline
    \rowcolor[gray]{0.8}
    \textbf{Competence} & \textbf{Description} & \textbf{Level of Mastery} \\ \hline
    Data Preprocessing & Cleaning, scaling, and transforming IoT data for machine learning tasks. & Advanced \\ \hline
    Time-Series Analysis & Classification and alignment techniques for temporal data from IoT sensors. & Intermediate \\ \hline
    Deep Learning for IoT & Designing and optimizing CNNs for resource-constrained environments. & Advanced \\ \hline
    Edge AI and Federated Learning & Implementing distributed learning paradigms while preserving privacy and improving performance. & Intermediate \\ \hline
    TinyML and Model Optimization & Applying pruning, quantization, and transfer learning to create lightweight AI models. & Intermediate \\ \hline
    Python Libraries & Proficiency with TensorFlow, TensorFlow Lite, scikit-learn, and Keras for embedded AI development. & Intermediate \\ \hline
    \end{tabular}
    \caption{Skills Acquired During the Embedded AI for IoT Module}
\end{table}

\subsection{Analysis and Remarks}
This course was very engaging as it allowed me to better understand some of the concepts I had previously encountered during a LinkedIn Learning course by Professor Céline Robardet from INSA Lyon. 
The practical sessions, especially the fall detection project, were particularly rewarding. 
It was gratifying to see how techniques like CNN optimization and federated learning could be applied to solve real-world IoT problems.


\newpage
% Emerging Networks
\section{Emerging Networks}
\subsection{Context}
\subsection{Technical Summary}
\subsection{Skills Acquired}
\subsection{Analysis and remarks}

\newpage
% Sensors Introduction
\section{Sensors Introduction}
\subsection{Context}
\subsection{Technical}
\subsection{Practical work}
\subsection{Skills acquired}
\subsection{Analysis and remarks}

\newpage
% Master REOC
\section{Master REOC}

\subsection{Context}
\indent \indent The Master REOC program is a 6 months program done in parallel of the 5ISS, the Thursdays afternoons. It is done in collaboration between INSA Toulouse and ENSEEIHT.

The objective of this program is to train network and telecommunications engineers, system/network architects, system/network administrators, and security engineers.

\subsection{Technical summary}
%Le Master REOC est structuré de manière à nous fournir à la fois des connaissances théoriques sur des problématiques spécifiques aux réseaux et télécommunications, et des
%compétences pratiques pour mettre en application les solutions disponibles pour les résoudre.
%Le Master se compose de 10 matières différentes, chacune abordant un aspect spécifique des réseaux et télécommunications.\\
    % Liste des cours
    %- Réseaux Embarqués et QoS
    %- Virtualisation dans l’embarqué
    %- Virtualisation de réseau / SDN
    %- Évaluation de performance Sans fil
    %- Virtualisation de réseau / NFV
    %- Virtualisation de réseau / Projet
    %- Architectures Modulaires
    %- Simulation avancée
    %- Pire cas avion
    %- Modélisation d’un Middleware IoT
    %- Modélisation d’un réseau LoRa"
%Nous avons abordés des sujets tels que la virtualisation, la qualité de service, la modélisation de réseaux, la simulation de réseaux, la sécurité des réseaux, la performance
%des réseaux sans fil, la modélisation de middleware IoT, la modélisation de réseaux LoRa,, mais aussi des sujets plus spécifiques comme les modules intégré de l'avionique ou l'analyse du temps de traversée pire cas sur les réseaux ethernet commutés.

\indent \indent The Master REOC is structured to provide both theoretical knowledge on specific networking and telecommunications challenges, as well as practical skills to implement available solutions to address these issues.

The program consists of 10 distinct courses, each focusing on a specific aspect of networking and telecommunications.
\\
We explored topics such as virtualization, quality of service (QoS), network modeling, network simulation, network security, wireless network performance, IoT middleware modeling, and LoRa network modeling.
Additionally, the program covered more specialized subjects, such as integrated avionics modules and worst-case traversal time analysis in switched Ethernet networks.

\subsection{Skills acquired}
Through this course, I consider to have acquired the following skills:

%- Conception de nouvelles solutions d'architectures réseau et de services associant l'ensemble des briques nécessaires (infrastructure, SI, réseau,....) en réponse à l'expression des besoins des opérateurs, entreprises, institutions privées ou publiques....
%- Réalisation d'une étude d'ingénierie détaillée afin de faire correspondre les déploiements locaux aux exigences de capacité, de couverture et de qualité de service définies dans le dossier d'architecture de communication.
%- Pilotage de l'implémentation des éléments de réseau et de l'intégration technique des équipements par les équipes opérationnelles suivant la nature des projets de déploiement.
%- Mise en service, paramétrage et configuration des équipements de réseaux, télécoms et services dans le cadre des installations prévues.
%- Respect du plan de prévention des risques et de l'application des règles de sécurité.
%- Supervision des systèmes de télécommunications, des équipements du réseau et des services au moyen des outils de supervision de son domaine.
%- Proposition, identification et définition des actions d'évolution et d'amélioration de service (à destination des équipes exploitation et/ou ingénierie).
%- Veille technologique et force de proposition sur de nouvelles fonctionnalités à ajouter aux solutions de services en développement.

% Embedded Networks and QoS
\begin{table}[H]
    \centering
    \arrayrulecolor{black} % Defines the color of the table borders
    \renewcommand{\arraystretch}{1.5} % Adjusts the vertical spacing between rows
    \begin{tabular}{|p{3.5cm}|p{8cm}|c|}
    \hline
    \rowcolor[gray]{0.8}
    \textbf{Competence} & \textbf{Description} & \textbf{Level of Mastery} \\
    \hline
        QoS Design & Designed networks with guaranteed Quality of Service for mission-critical applications. & Advanced \\ \hline
        Network Constraint Analysis & Evaluated latency, throughput, and reliability constraints in embedded systems. & Intermediate \\ \hline
        \end{tabular}
        \caption{Skills Acquired in the "Embedded Networks and QoS" course}
        \end{table}
        
% Virtualization in Embedded Systems
\begin{table}[H]
    \centering
    \arrayrulecolor{black} % Defines the color of the table borders
    \renewcommand{\arraystretch}{1.5} % Adjusts the vertical spacing between rows
    \begin{tabular}{|p{3.5cm}|p{8cm}|c|}
        \hline
        \rowcolor[gray]{0.8}
        \textbf{Competence} & \textbf{Description} & \textbf{Level of Mastery} \\
        \hline
        Embedded System Virtualization & Implemented virtualization techniques to enhance resource utilization in embedded devices. & Advanced \\ \hline
        System Partitioning & Applied partitioning principles to improve modularity and scalability. & Intermediate \\ \hline
    \end{tabular}
\caption{Skills Acquired in the "Virtualization in Embedded Systems"course}
\end{table}
    
% Wireless Network Performance Evaluation
\begin{table}[H]
    \centering
    \arrayrulecolor{black} % Defines the color of the table borders
    \renewcommand{\arraystretch}{1.5} % Adjusts the vertical spacing between rows
    \begin{tabular}{|p{3.5cm}|p{8cm}|c|}
        \hline
        \rowcolor[gray]{0.8}
        \textbf{Competence} & \textbf{Description} & \textbf{Level of Mastery} \\
        \hline
        Wireless Network Optimization & Analyzed and optimized performance metrics like latency and throughput in wireless systems. & Advanced \\ \hline
        Real-Time Performance Evaluation & Conducted real-time performance analysis for wireless networks in constrained environments. & Intermediate \\ \hline
    \end{tabular}
\caption{Skills Acquired in the "Wireless Network Performance Evaluation" course}
\end{table}
    
% Modular Architectures
\begin{table}[H]
    \centering
    \arrayrulecolor{black} % Defines the color of the table borders
    \renewcommand{\arraystretch}{1.5} % Adjusts the vertical spacing between rows
    \begin{tabular}{|p{3.5cm}|p{8cm}|c|}
        \hline
        \rowcolor[gray]{0.8}
        \textbf{Competence} & \textbf{Description} & \textbf{Level of Mastery} \\
        \hline
        ARINC 653 Principles & Applied modular design principles to Integrated Modular Avionics (IMA). & Advanced \\ \hline
        Partitioning for Isolation & Designed systems with spatial and temporal isolation using ARINC 653. & Intermediate \\ \hline
    \end{tabular}
\caption{Skills Acquired in the "Modular Architectures" course}
\end{table}
    
% Worst-Case Traversal Time Analysis
\begin{table}[H]
    \centering
    \arrayrulecolor{black} % Defines the color of the table borders
    \renewcommand{\arraystretch}{1.5} % Adjusts the vertical spacing between rows
    \begin{tabular}{|p{3.5cm}|p{8cm}|c|}
    \hline
    \rowcolor[gray]{0.8}
    \textbf{Competence} & \textbf{Description} & \textbf{Level of Mastery} \\
    \hline
        WCTT Calculation & Computed worst-case traversal time using network calculus for real-time Ethernet. & Intermediate \\ \hline
        Latency Optimization & Optimized network latency to meet real-time constraints. & Advanced \\ \hline
        \end{tabular}
        \caption{Skills Acquired in the "Worst-Case Traversal Time Analysis" course}
        \end{table}
        
        % IoT Middleware and LoRa Modeling
\begin{table}[H]
    \centering
    \arrayrulecolor{black} % Defines the color of the table borders
    \renewcommand{\arraystretch}{1.5} % Adjusts the vertical spacing between rows
    \begin{tabular}{|p{3.5cm}|p{8cm}|c|}
        \hline
        \rowcolor[gray]{0.8}
        \textbf{Competence} & \textbf{Description} & \textbf{Level of Mastery} \\
        \hline
        IoT Middleware Design & Modeled middleware for efficient communication in IoT systems. & Advanced \\ \hline
        LoRa Network Simulation & Simulated and evaluated LoRa networks for low-power IoT applications. & Advanced \\ \hline
    \end{tabular}
    \caption{Skills Acquired in the "IoT Middleware and LoRa Modeling" course}
\end{table}

\begin{table}[H]
    \centering
    \arrayrulecolor{black} % Defines the color of the table borders
    \renewcommand{\arraystretch}{1.5} % Adjusts the vertical spacing between rows
    \begin{tabular}{|p{3.5cm}|p{8cm}|c|}
    \hline
    \rowcolor[gray]{0.8}
    \textbf{Competence} & \textbf{Description} & \textbf{Level of Mastery} \\
    \hline
    Defining and Deploying Virtual Networks & Designing and implementing virtual networks, ensuring proper communication between components. & Advanced \\
    \hline
    Working with SDN and NFV & Gained hands-on experience creating VNFs, managing SDN controllers, and addressing network challenges. & Advanced \\
    \hline
    Network Monitoring and Optimization & Implemented systems to track bandwidth and mitigate congestion through SDN rule adjustments. & Advanced\\
    \hline
    Collaborative Problem-Solving & Successfully worked in a team to design and implement a complex network simulation. & Advanced\\
    \hline
    \end{tabular}
    \caption{Competences Gained in the "Network Virtualization" (SDN \& NFV) Course}
\end{table}

\subsection{Analysis and remarks}
\indent \indent The Master REOC program was one of the most demanding experiences of my academic journey.
The practical sessions, particularly the Mininet project, required intense focus and teamwork to complete within the constraints of the lab.
While the workload was challenging, it was incredibly rewarding to see our theoretical knowledge applied to real-world scenarios.

This project helped me develop a deeper understanding of constrained networks and cutting-edge technologies like SDN and NFV.
The process of designing and simulating a network topology taught me the importance of planning and precision, while implementing VNFs and developing the Flask interface allowed me to enhance my technical and problem-solving skills.

If I had the opportunity to redo this year, I would organise myself better.
The challenges and time constraints, in parallel of my courses at INSA, were quite imposing, and pushed me to improve both my technical abilities and my ability to work under pressure.
In fact, I think that this was a great experience, but I would have liked to have more time to explore the different courses.

\newpage
\chapter{Conclusion}
\end{document}
